% ============================================================
%  EXAMEN TEMA 4: CINEMÁTICA — MRU, MRUV Y CAÍDA LIBRE
%  Curso 4to Año — Física
%  Duración: 90 minutos
%  Temas: Movimiento rectilíneo uniforme, uniformemente variado,
%         caída libre y lanzamiento vertical
%  Autor: Ing. Gabriel Astudillo
%  Nota: enunciado limpio (sin pistas ni desarrollo). Las
%  soluciones están en la "Clave de Respuestas" al final.
% ============================================================

\documentclass[12pt, letterpaper]{article}

% ── Paquetes ─────────────────────────────────────────────────

\usepackage{mathwizards-guia}
\mwguia{Examen — Tema 4: Cinemática}{Ing. Gabriel Astudillo $\cdot$ 4to Año}

\begin{document}
% ============================================================

% ── Portada / Título ─────────────────────────────────────────
\mwportada{Examen del Tema 4}{Cinemática: el estudio del movimiento}{MRU, MRUV, Caída Libre y Lanzamiento Vertical}

\vspace{4pt}
\begin{cuadroinfo}
  \small
  \textbf{Nombre:} \rule{0.55\linewidth}{0.4pt} \quad \textbf{Fecha:} \rule{0.15\linewidth}{0.4pt} \\[6pt]
  \textbf{Tiempo disponible:} 90 minutos \quad$\bullet$\quad \textbf{Puntaje total:} 100 puntos \\[4pt]
  \textbf{Instrucciones:}
  \begin{enumerate}[leftmargin=*]
    \item Lee cada ejercicio con calma antes de resolverlo.
    \item Muestra \emph{todos} tus pasos: datos, fórmula, sustitución y respuesta.
    \item \textbf{No olvides las unidades} y convierte cuando sea necesario.
    \item Usa $g = 9{,}8\;\text{m/s}^2$ salvo que se indique otra cosa.
    \item No se permite el uso de calculadora.
    \item Escribe con letra clara y revisa tu examen antes de entregarlo.
  \end{enumerate}
\end{cuadroinfo}

\vspace{8pt}

% ============================================================
\section{Movimiento Rectilíneo Uniforme (MRU) (30 puntos)}
% ============================================================

\begin{enumerate}[leftmargin=*]
  \item \textbf{(10 pts)} Resuelve:
        \begin{enumerate}[label=\alph*)]
          \item Un auto parte de la posición $x_0 = 20\;\text{km}$ y viaja a $72\;\text{km/h}$ durante $2{,}5\;\text{h}$. \textquestiondown Cuál es su posición final?
          \item Un tren recorre $270\;\text{km}$ en $2{,}25\;\text{h}$. \textquestiondown Cuál es su rapidez media?
          \item \textquestiondown Qué distancia recorre un móvil a $15\;\text{m/s}$ en $12\;\text{s}$?
        \end{enumerate}

  \item \textbf{(10 pts)} Dos ciudades están separadas por $450\;\text{km}$. Dos autos parten simultáneamente uno hacia el otro: el primero a $70\;\text{km/h}$ y el segundo a $80\;\text{km/h}$. \textquestiondown Después de cuántas horas se encuentran y a qué distancia de la primera ciudad?

  \item \textbf{(10 pts)} Un auto $A$ sale de un punto a $60\;\text{km/h}$. Una hora después, el auto $B$ sale del mismo punto a $90\;\text{km/h}$ en la misma dirección. \textquestiondown Después de cuántas horas (desde que sale $B$) lo alcanza y a qué distancia del punto de partida?
\end{enumerate}

\newpage

% ============================================================
\section{Movimiento Rectilíneo Uniformemente Variado (MRUV) (35 puntos)}
% ============================================================

\begin{enumerate}[leftmargin=*]
  \item \textbf{(9 pts)} Calcula:
        \begin{enumerate}[label=\alph*)]
          \item La aceleración de un móvil que pasa de $8\;\text{m/s}$ a $20\;\text{m/s}$ en $4\;\text{s}$.
          \item La velocidad que alcanza un auto que parte del reposo con $a = 5\;\text{m/s}^2$ durante $6\;\text{s}$.
          \item La distancia de frenado de un auto a $30\;\text{m/s}$ que frena con $a = -6\;\text{m/s}^2$ hasta detenerse.
        \end{enumerate}

  \item \textbf{(8 pts)} Un tren parte del reposo con $a = 3\;\text{m/s}^2$ durante $10\;\text{s}$ y luego mantiene su velocidad constante durante $15\;\text{s}$. \textquestiondown Qué distancia total recorre?

  \item \textbf{(10 pts)} Un auto que viaja a $25\;\text{m/s}$ frena uniformemente con $a = -5\;\text{m/s}^2$ hasta detenerse.
        \begin{enumerate}[label=\alph*)]
          \item \textquestiondown Qué distancia recorre hasta detenerse?
          \item \textquestiondown Cuánto tiempo tarda en detenerse?
        \end{enumerate}

  \item \textbf{(8 pts)} Un móvil parte con $v_i = 12\;\text{m/s}$ y acelera a $2\;\text{m/s}^2$.
        \begin{enumerate}[label=\alph*)]
          \item \textquestiondown Qué velocidad tiene a los $8\;\text{s}$?
          \item \textquestiondown Qué distancia recorre en esos $8\;\text{s}$?
        \end{enumerate}
\end{enumerate}

\newpage

% ============================================================
\section{Caída Libre y Lanzamiento Vertical (35 puntos)}
% ============================================================

\begin{enumerate}[leftmargin=*]
  \item \textbf{(9 pts)} Una piedra se deja caer desde un puente y tarda $4\;\text{s}$ en llegar al agua ($g = 9{,}8\;\text{m/s}^2$).
        \begin{enumerate}[label=\alph*)]
          \item \textquestiondown Qué altura tiene el puente?
          \item \textquestiondown Con qué velocidad llega la piedra al agua?
        \end{enumerate}

  \item \textbf{(9 pts)} Una pelota se lanza verticalmente hacia arriba con $v_i = 29{,}4\;\text{m/s}$.
        \begin{enumerate}[label=\alph*)]
          \item \textquestiondown Cuál es la altura máxima que alcanza?
          \item \textquestiondown Cuál es el tiempo total de vuelo (hasta volver al punto de lanzamiento)?
        \end{enumerate}

  \item \textbf{(8 pts)} Desde un edificio de $45\;\text{m}$ de altura se lanza una pelota verticalmente \emph{hacia abajo} con $v_i = 5\;\text{m/s}$. \textquestiondown Con qué velocidad llega al suelo? (Usa $v^2 = v_i^2 + 2gh$.)

  \item \textbf{(9 pts)} Una pelota se lanza hacia arriba con $v_i = 19{,}6\;\text{m/s}$. Usando $h = v_i t - \frac{1}{2} g t^2$, determina los \textbf{instantes} en que la pelota se encuentra a $14{,}7\;\text{m}$ de altura. \textquestiondown Qué relación tienen esos dos instantes?
\end{enumerate}

\newpage

% ============================================================
\ifmwclaves
  \section{Clave de Respuestas}
  % ============================================================

  \subsection*{Sección 1: Movimiento Rectilíneo Uniforme (MRU)}

  \begin{enumerate}[leftmargin=*, label=\textbf{\arabic*.}]
    \item a) $x = x_0 + v\,t = 20 + 72(2{,}5) = 20 + 180 = 200\;\text{km}$.
          \quad b) $v = \dfrac{270}{2{,}25} = 120\;\text{km/h}$.
          \quad c) $d = v\,t = 15 \times 12 = 180\;\text{m}$.

    \item Velocidad de acercamiento: $70 + 80 = 150\;\text{km/h}$. \quad $t = \dfrac{450}{150} = 3\;\text{h}$.
          Distancia desde la primera ciudad: $70 \times 3 = 210\;\text{km}$.

    \item Cuando $B$ sale, $A$ lleva una ventaja de $60 \times 1 = 60\;\text{km}$.
          Velocidad de acercamiento: $90 - 60 = 30\;\text{km/h}$. \quad $t = \dfrac{60}{30} = 2\;\text{h}$ (desde que sale $B$).
          Distancia del punto de partida: $90 \times 2 = 180\;\text{km}$.
  \end{enumerate}

  \newpage

  \subsection*{Sección 2: Movimiento Rectilíneo Uniformemente Variado (MRUV)}

  \begin{enumerate}[leftmargin=*, label=\textbf{\arabic*.}]
    \item a) $a = \dfrac{20 - 8}{4} = 3\;\text{m/s}^2$.
          \quad b) $v_f = v_i + a\,t = 0 + 5(6) = 30\;\text{m/s}$.
          \quad c) $v_f^2 = v_i^2 + 2a\,\Delta x \Rightarrow 0 = 30^2 + 2(-6)\Delta x \Rightarrow 12\,\Delta x = 900 \Rightarrow \Delta x = 75\;\text{m}$.

    \item Tramo 1 (acelerado, $10\;\text{s}$): $v_f = 3(10) = 30\;\text{m/s}$; $\Delta x_1 = \tfrac{1}{2}(3)(10)^2 = 150\;\text{m}$.
          \\[2pt]
          Tramo 2 (MRU, $15\;\text{s}$): $\Delta x_2 = 30 \times 15 = 450\;\text{m}$.
          \\[2pt]
          Distancia total: $150 + 450 = 600\;\text{m}$.

    \item a) $v_f^2 = v_i^2 + 2a\,\Delta x \Rightarrow 0 = 25^2 + 2(-5)\Delta x \Rightarrow 10\,\Delta x = 625 \Rightarrow \Delta x = 62{,}5\;\text{m}$.
          \quad b) $v_f = v_i + a\,t \Rightarrow 0 = 25 - 5t \Rightarrow t = 5\;\text{s}$.

    \item a) $v_f = 12 + 2(8) = 28\;\text{m/s}$.
          \quad b) $\Delta x = v_i t + \tfrac{1}{2}a\,t^2 = 12(8) + \tfrac{1}{2}(2)(8)^2 = 96 + 64 = 160\;\text{m}$.
  \end{enumerate}

  \newpage

  \subsection*{Sección 3: Caída Libre y Lanzamiento Vertical}

  \begin{enumerate}[leftmargin=*, label=\textbf{\arabic*.}]
    \item a) $h = \tfrac{1}{2}g\,t^2 = \tfrac{1}{2}(9{,}8)(4)^2 = 4{,}9 \times 16 = 78{,}4\;\text{m}$.
          \quad b) $v = g\,t = 9{,}8 \times 4 = 39{,}2\;\text{m/s}$.

    \item a) $h_{\max} = \dfrac{v_i^2}{2g} = \dfrac{(29{,}4)^2}{2(9{,}8)} = \dfrac{864{,}36}{19{,}6} = 44{,}1\;\text{m}$.
          \quad b) $t_{\text{vuelo}} = \dfrac{2v_i}{g} = \dfrac{2(29{,}4)}{9{,}8} = 6\;\text{s}$.

    \item $v^2 = v_i^2 + 2gh = 5^2 + 2(9{,}8)(45) = 25 + 882 = 907$; \quad $v = \sqrt{907} \approx 30{,}1\;\text{m/s}$.

    \item Planteamos $19{,}6\,t - 4{,}9\,t^2 = 14{,}7$:
          \[
            4{,}9\,t^2 - 19{,}6\,t + 14{,}7 = 0 \Rightarrow t = \frac{19{,}6 \pm \sqrt{384{,}16 - 288{,}12}}{9{,}8} = \frac{19{,}6 \pm 9{,}8}{9{,}8}.
          \]
          $t_1 = 1\;\text{s}$ y $t_2 = 3\;\text{s}$.
          Los dos instantes son simétricos respecto al tiempo de altura máxima ($t = 2\;\text{s}$): la pelota pasa por esa altura subiendo ($1\;\text{s}$) y bajando ($3\;\text{s}$).
  \end{enumerate}

\fi
\end{document}
